\chapter{परिमित विमा}\label{ch:b1:findim}

परिमित कितने सदिशों से \omterm{def:b1:vspaces:span}{जनित समष्टि} एक सुपरिभाषित \emph{विमा} धारण करती है
--- अर्थात् उसके सभी आधारों का उभयनिष्ठ आकार। नीचे की सारी उपपत्तियाँ एक ही
संचयात्मक यंत्र से बहती हैं, विनिमय प्रमेयिका: कोई \omterm{def:b1:vspaces:free}{स्वतंत्र कुल} किसी जनक कुल
से बड़ा कभी नहीं हो सकता। विमा के साथ वे औज़ार आते हैं जो आगे सर्वत्र काम
आते हैं: अपूर्ण \omterm{def:b1:vspaces:free}{आधार} प्रमेय, किसी कुल की \omterm{def:b1:findim:rank}{कोटि}, ग्रासमान का सूत्र।

\section{आधारों का अस्तित्व}

\begin{definition}\label{def:b1:findim:def}
$E$ \emph{परिमित-विमीय}\index{परिमित विमा} तब कहलाती है जब उसका कोई परिमित
जनक कुल हो। (अन्यथा \emph{अनंत-विमीय}: जैसे $K[X]$, जिसके परिमित कुल केवल
परिबद्ध घात वाले \omterm{def:b1:poly:def}{बहुपद} जनित करते हैं।)
\end{definition}

\begin{theorem}[विनिमय प्रमेयिका]\label{thm:b1:findim:exchange}
मान लीजिए $(g_1, \dots, g_n)$ $E$ को जनित करता है और $(f_1, \dots, f_p)$ $E$
में \omterm{def:b1:vspaces:free}{स्वतंत्र} है। तब $p \leq n$।
\end{theorem}

\begin{proof}
हम $k \leq p$ पर आगमन से सिद्ध करते हैं: \emph{$g_j$ को फिर से क्रमांकित
करने के बाद कुल $(f_1, \dots, f_k, g_{k+1}, \dots, g_n)$ $E$ को जनित करता
है} --- जो हर चरण पर $k \leq n$ बाध्य कर देता है, और अंत में $p \leq n$।

$k = 0$: यही परिकल्पना है। चरण: उसे $k - 1$ के लिए मान लीजिए; तब $f_k$
$(f_1, \dots, f_{k-1},$ $g_k, \dots, g_n)$ का संयोजन है। इस संयोजन में किसी
$g_j$ ($j \geq k$) का गुणांक अशून्य है --- अन्यथा $f_k$ $f_1, \dots,
f_{k-1}$ का संयोजन होता, जो स्वतंत्रता का विरोध करता। इस प्रकार फिर से
क्रमांकित कीजिए कि $j = k$, और $g_k$ के लिए हल कीजिए: $g_k$ $(f_1, \dots,
f_k, g_{k+1}, \dots, g_n)$ का संयोजन है। तब $E$ का हर सदिश, जो $(k-1)$-कुल
के द्वारा व्यक्त था, $k$-कुल के द्वारा फिर से व्यक्त किया जा सकता है: अतः वह
जनक है। (यदि $k - 1 = n$, तो कोई $g$ शेष नहीं रहता और $f_k$ अकेले $f_i$ का
संयोजन होता: असंभव; अतः $k \leq n$।)
\end{proof}

\begin{example}[विनिमय, एक बार देखकर]
\label{ex:b1:findim:exchangerun}
$\R^2$ में जनक कुल $\bigl(g_1, g_2\bigr) = \bigl((1,0), (0,1)\bigr)$ और
\omterm{def:b1:vspaces:free}{स्वतंत्र कुल} $\bigl(f_1, f_2\bigr) = \bigl((1,2), (3,4)\bigr)$ लीजिए। चरण
$1$: $f_1 = 1\cdot g_1 + 2\cdot g_2$; $g_2$ का गुणांक अशून्य है, अतः $g_2$
के बदले $f_1$ लाइए: कुल $\bigl(f_1, g_1\bigr)$ अब भी जनक है ($g_2 =
\frac12(f_1 - g_1)$)। चरण $2$: $f_2 = (3,4) = 2\,f_1 + 1\cdot g_1$; शेष
$g_1$ का गुणांक अशून्य है (होना ही है: $f_2$ $f_1$ का गुणज नहीं है), अतः फिर
से विनिमय कीजिए: $\bigl(f_1, f_2\bigr)$ $\R^2$ को जनित करता है। यदि कोई
तीसरा \omterm{def:b1:vspaces:free}{स्वतंत्र} सदिश $f_3$ होता, तो उसे सोखने के लिए कोई $g$ बचता ही नहीं ---
और ठीक इसी तरह यह प्रमेयिका $\R^2$ में $3$ \omterm{def:b1:vspaces:free}{स्वतंत्र} सदिशों को मना करती है।
ऊपर की उपपत्ति यही लेखा-जोखा व्यापक रूप में किया हुआ है।
\end{example}

\begin{theorem}[परिमित विमा में आधार]\label{thm:b1:findim:bases}
मान लीजिए $E \neq \{0\}$ \omterm{def:b1:findim:def}{परिमित-विमीय} है।
\begin{enumerate}
  \item हर परिमित जनक कुल में से एक \omterm{def:b1:vspaces:free}{आधार} निकाला जा सकता है।
  \item (अपूर्ण \omterm{def:b1:vspaces:free}{आधार} प्रमेय\index{अपूर्ण आधार प्रमेय}) हर \omterm{def:b1:vspaces:free}{स्वतंत्र कुल} को
  \omterm{def:b1:vspaces:free}{आधार} तक बढ़ाया जा सकता है, और वह भी किसी भी चुने हुए जनक कुल के सदिशों से।
  \item $E$ के सभी \omterm{def:b1:vspaces:free}{आधार} परिमित हैं और उनमें अवयवों की संख्या समान है: यही
  \emph{विमा} $\dim E$ है। (परिपाटी: $\dim\{0\} = 0$।)
\end{enumerate}
\end{theorem}

\begin{proof}
(1) एक-एक करके ऐसा हर सदिश हटा दीजिए जो शेष का संयोजन हो; कुल जनक बना रहता
है, क्योंकि हटाए गए सदिश का प्रयोग करने वाले किसी भी व्यंजक में बचे हुओं का
उसका संयोजन रख दिया जा सकता है। यह प्रक्रिया समाप्त होती है --- हर चरण
परिमित कुल को एक से छोटा करता है --- और ठीक तब रुकती है जब शेष कोई भी सदिश
बाकी सदिशों का संयोजन न हो। अंतिम कुल तब भी जनक है, और \omterm{def:b1:vspaces:free}{स्वतंत्र} भी: किसी
अतुच्छ शून्य संयोजन में कोई अशून्य गुणांक होता, और उससे भाग देने पर उस सदिश
को शेष के पदों में हल किया जा सकता, जिससे वह हटाने योग्य निकलता --- और यह इस
बात का विरोध करता कि प्रक्रिया रुक चुकी थी।

(2) मान लीजिए $(f_1, \dots, f_p)$ \omterm{def:b1:vspaces:free}{स्वतंत्र} है और $(g_1, \dots, g_n)$ जनक।
$g_1, \dots, g_n$ पर चलते जाइए, और जब भी $g_j$ वर्तमान कुल की \omterm{def:b1:vspaces:span}{जनित समष्टि}
में पहले से न हो, उसे कुल में जोड़ दीजिए
(\cref{prop:b1:vspaces:freecriteria}~(2) कुल को \omterm{def:b1:vspaces:free}{स्वतंत्र} बनाए रखता है)।
अंतिम कुल \omterm{def:b1:vspaces:free}{स्वतंत्र} है, और जनक भी: हर $g_j$ उसकी \omterm{def:b1:vspaces:span}{जनित समष्टि} में है --- या तो
वह जोड़ा गया था, या वह पहले से ही एक संयोजन था।

(3) दो आधारों में से प्रत्येक \omterm{def:b1:vspaces:free}{स्वतंत्र} भी है और जनक भी: विनिमय प्रमेयिका
उनकी गणनसंख्याओं के बीच दोनों असमिकाएँ दे देती है। (परिमितता: \omterm{def:b1:vspaces:free}{आधार} \omterm{def:b1:vspaces:free}{स्वतंत्र}
होता है, अतः विनिमय से वह किसी परिमित जनक कुल से बड़ा नहीं हो सकता।)
\end{proof}

\begin{example}\label{ex:b1:findim:dims}
$\dim K^n = n$ (विहित \omterm{def:b1:vspaces:free}{आधार}); $\dim K_n[X] = n + 1$ (एकपदी); $\dim_\R \C =
2$; और $y'' + ay' + by = 0$ की हल-समष्टि की विमा $2$ है
(\cref{thm:b1:diffeq:homogeneous2}: हल $(\lambda, \mu) \in K^2$ के द्वारा
\omterm{def:b1:logic:inj}{एकैकी आच्छादक} तथा रैखिक रूप से प्राचलित होते हैं)।
\end{example}

\begin{example}[एक ही समुच्चय, दो विमाएँ]
\label{ex:b1:findim:twodimensions}
सम्मिश्र संख्याओं के युग्मों का \omterm{def:b1:logic:sets}{समुच्चय} $\C^2$ विमा $2$ वाली $\C$-सदिश
समष्टि है (विहित \omterm{def:b1:vspaces:free}{आधार} $e_1, e_2$) --- और विमा $4$ वाली $\R$-सदिश समष्टि भी,
जिसका \omterm{def:b1:vspaces:free}{आधार} है
\[
(1, 0),\quad (\iu, 0),\quad (0, 1),\quad (0, \iu):
\]
हर $(z, w) = (a + \iu b,\ c + \iu d)$ के वास्तविक \omterm{prop:b1:vspaces:coordinates}{निर्देशांक} $(a, b, c, d)$
अद्वितीय रूप से होते हैं। विमा केवल सदिशों के \omterm{def:b1:logic:sets}{समुच्चय} का गुण नहीं है: वह
\emph{अनुमत अदिशों के सापेक्ष} स्वातंत्र्य कोटियाँ गिनती है, और अदिशों की
पूर्ति $\C$ से $\R$ तक आधी कर देने पर गिनती दुगुनी हो जाती है। (सप्ताहांत
समस्या इसी संवेदनशीलता की चरम स्थिति का लाभ उठाती है, जहाँ अदिश सिकुड़कर
$\Q$ तक आ जाते हैं।)
\end{example}

\begin{example}[पूर्णन कलनविधि चलाकर देखना]
\label{ex:b1:findim:completion}
विहित सदिशों का प्रयोग करके \omterm{def:b1:vspaces:free}{स्वतंत्र कुल} $\bigl((1,1,1)\bigr)$ को $\R^3$ के
\omterm{def:b1:vspaces:free}{आधार} तक पूर्ण कीजिए। जनक कुल $(e_1, e_2, e_3)$ पर
\cref{thm:b1:findim:bases}~(2) की उपपत्ति चलाइए: क्या $e_1 \in
\operatorname{Vect}(1,1,1)$? नहीं ($(1,1,1)$ के गुणजों के \omterm{prop:b1:vspaces:coordinates}{निर्देशांक} समान
होते हैं) --- अतः उसे जोड़ दीजिए। क्या $e_2 \in
\operatorname{Vect}\bigl((1,1,1), e_1\bigr)$? किसी संयोजन $\alpha(1,1,1) +
\beta e_1$ के दूसरे और तीसरे \omterm{prop:b1:vspaces:coordinates}{निर्देशांक} समान होते हैं, और $e_2$ के नहीं ---
अतः उसे भी जोड़ दीजिए। कुल $\bigl((1,1,1), e_1, e_2\bigr)$ $3$ सदिशों के साथ
\omterm{def:b1:vspaces:free}{स्वतंत्र} है: रुक जाइए, वही \omterm{def:b1:vspaces:free}{आधार} है (\cref{prop:b1:findim:twoofthree} इस
प्रतिवर्त को औपचारिक बना देगा)। ध्यान दीजिए कि परिणाम इस बात पर निर्भर करता
है कि $g_j$ किस क्रम में देखे गए: पूर्णन एक कलनविधि है, कोई सूत्र नहीं।
\end{example}

\begin{proposition}[तीन में से दो वाला नियम]\label{prop:b1:findim:twoofthree}
मान लीजिए $\dim E = n$ और $\mathcal{F}$ $E$ के ठीक $n$ सदिशों का कुल है। तब
\[
\mathcal{F} \text{ स्वतंत्र} \iff \mathcal{F} \text{ जनक} \iff
\mathcal{F} \text{ आधार}.
\]
इसके अतिरिक्त किसी भी \omterm{def:b1:vspaces:free}{स्वतंत्र कुल} में $\leq n$ सदिश होते हैं, और किसी भी
जनक कुल में $\geq n$।
\end{proposition}

\begin{proof}
गणनसंख्या के परिबंध किसी \omterm{def:b1:vspaces:free}{आधार} के सामने लगाई गई विनिमय प्रमेयिका ही हैं। यदि
$\mathcal{F}$ (आकार $n$) \omterm{def:b1:vspaces:free}{स्वतंत्र} है पर जनक नहीं, तो कोई $x$ उसकी \omterm{def:b1:vspaces:span}{जनित
समष्टि} के बाहर है; $x$ जोड़ने पर $n + 1$ सदिशों का \omterm{def:b1:vspaces:free}{स्वतंत्र कुल} मिल जाता:
असंभव। यदि $\mathcal{F}$ जनक है पर \omterm{def:b1:vspaces:free}{स्वतंत्र} नहीं, तो \omterm{def:b1:vspaces:free}{आधार} निकालने पर
(\cref{thm:b1:findim:bases}~(1)) $< n$ सदिशों वाला \omterm{def:b1:vspaces:free}{आधार} मिल जाता: असंभव।
\end{proof}

\begin{example}[तीन में से दो, आधा काम बचाकर]
\label{ex:b1:findim:twoofthreerun}
क्या $\bigl((1,1,0), (0,1,1), (1,0,1)\bigr)$ $\R^3$ का \omterm{def:b1:vspaces:free}{आधार} है? गिनती: तीन
सदिश, विमा तीन --- अतः निर्णय अकेली स्वतंत्रता कर देगी। शून्य संयोजन से $a +
c = 0$, $a + b = 0$, $b + c = 0$ मिलता है; तीनों को जोड़ने पर $2(a + b + c)
= 0$, और $a + b + c = 0$ में से हर मूल समीकरण घटाने पर $b = c = a = 0$ बचता
है: अतः \omterm{def:b1:vspaces:free}{स्वतंत्र}, इसलिए \omterm{def:b1:vspaces:free}{आधार}, और सत्यापन का जनक-वाला आधा हिस्सा प्रमेय ने
मुफ़्त में दे दिया। \cref{ex:b1:vspaces:testbasis} से तुलना कीजिए, जहाँ वही
दुहरा सत्यापन हाथ से करना पड़ा था --- सिद्धांत का एक अध्याय ठीक इतनी ही बचत
में बदल जाता है, और वह भी पुस्तक के शेष भाग की हर आधार-जाँच पर।
\end{example}

\begin{method}[विमा की संगणना]
\label{met:b1:findim:computedim}
तीन मानक रास्ते, घटती हुई बारंबारता के क्रम में।
\begin{enumerate}
  \item \emph{प्राचलन कीजिए, फिर \omterm{def:b1:vspaces:free}{आधार} पढ़ लीजिए।} परिभाषक प्रतिबंध हल कीजिए,
  व्यापक अवयव को बचे हुए प्राचलों में रैखिक रूप से व्यक्त कीजिए, और जाँचिए
  कि प्राचलों से गुणा होने वाले सदिश \omterm{def:b1:vspaces:free}{स्वतंत्र} हैं: विमा प्राचलों की संख्या
  है। (नीचे $\R^4$ की एक ठोस \omterm{def:b1:vspaces:subspace}{उपसमष्टि} पर चलाया गया है।)
  \item \emph{कोई \omterm{def:b1:logic:inj}{एकैकी आच्छादक} रैखिक प्राचलन दिखाइए।} जब अवयव परिमित कितने
  मानों से तय हो जाएँ --- किसी पुनरावृत्ति की आरंभिक शर्तें
  (\cref{exo:b1:findim:10}), किसी हल-सूत्र के गुणांक
  (\cref{ex:b1:findim:dims}) --- तब विमा उन मानों की संख्या है।
  \item \emph{सूत्रों का प्रयोग कीजिए।} सर्वनिष्ठ और योग के लिए ग्रासमान,
  जनित समष्टियों के लिए \omterm{def:b1:findim:rank}{कोटि}, और आगे चलकर अष्टि तथा प्रतिबिंब के लिए
  कोटि--शून्यता: विमाएँ प्रायः \emph{संगणित} की जाती हैं, अनुमानित नहीं।
\end{enumerate}
तीनों रास्तों में तीन में से दो वाला नियम ही समापक है: गिनती मिलते ही अकेली
स्वतंत्रता या अकेला जनक होना निष्कर्ष दे देता है।
\end{method}

\begin{example}[रास्ता 1, पूरा]
\label{ex:b1:findim:parametrize}
$H = \{(x, y, z, t) \in \R^4 : x + y + z + t = 0 \text{ और } x = t\}$ की
विमा। हल कीजिए: $t = x$ और $y + z = -2x$, अतः $x, y$ के मुक्त रहते $z = -2x
- y$:
\[
(x,\ y,\ -2x - y,\ x) = x\,(1, 0, -2, 1) + y\,(0, 1, -1, 0) .
\]
दोनों सदिश \omterm{def:b1:vspaces:free}{स्वतंत्र} हैं (पहले दो \omterm{prop:b1:vspaces:coordinates}{निर्देशांक} देखिए: $(x, y) = (0,0)$), अतः वे
$H$ का \omterm{def:b1:vspaces:free}{आधार} बनाते हैं और $\dim H = 2$। गिनती पहले से बताई जा सकती थी ---
$\R^4$ में दो \omterm{def:b1:vspaces:free}{स्वतंत्र} रैखिक प्रतिबंधों में से हर एक को एक विमा खा जानी
चाहिए --- पर प्राचलन उसे सिद्ध \emph{भी} करता है और एक \omterm{def:b1:vspaces:free}{आधार} भी थमा देता है,
जो अकेली भविष्यवाणी कभी नहीं करती; \cref{exo:b1:findim:9} ``हर समीकरण अधिक
से अधिक एक विमा खाता है'' वाले नारे को प्रमेय बना देता है।
\end{example}

\begin{example}[रास्ता 2, पूरा]
\label{ex:b1:findim:route2}
$W = \{P \in \R_n[X] : P(1) = P(2) = 0\}$ की विमा ($n \geq 2$ के लिए)।
गुणनखंड प्रमेय को दो बार लगाने पर (\cref{thm:b1:poly:factor}; मूल $1$ और $2$
भिन्न हैं), $P \in W$ ठीक तब होता है जब $\deg Q \leq n - 2$ के साथ $P = (X -
1)(X - 2)\,Q$। संगति $Q \mapsto (X-1)(X-2)Q$ रैखिक है, पूरे $W$ तक पहुँचती
है, और \omterm{def:b1:logic:inj}{एकैकी} है (गुणनफल शून्य तभी होता है जब $Q = 0$): अतः $W$ $\R_{n-2}[X]$
के द्वारा \omterm{def:b1:logic:inj}{एकैकी आच्छादक} तथा रैखिक रूप से प्राचलित है, इसलिए
\[
\dim W = \dim \R_{n-2}[X] = n - 1 .
\]
\omterm{def:b1:vspaces:free}{आधार} प्राचलन के साथ ही मिल जाता है: एकपदियों के प्रतिबिंब,
$\bigl((X-1)(X-2),\ (X-1)(X-2)X,\ \dots,\ (X-1)(X-2)X^{n-2}\bigr)$। किसी नए
बिंदु पर हर नया मूल्यांकन-प्रतिबंध ठीक एक विमा की क़ीमत लेता है --- यही
\omterm{thm:b1:poly:lagrange}{लाग्रांज अंतर्वेशन} (\cref{thm:b1:poly:lagrange}) की गिनती वाली रीढ़ है।
\end{example}

\begin{example}[समाकल वाला प्रतिबंध भी एक विमा लेता है]
\label{ex:b1:findim:integralconstraint}
$H = \{P \in \R_2[X] : \int_0^1 P = 0\}$ की विमा। $P = a + bX + cX^2$ लिखने
पर प्रतिबंध $a + \frac b2 + \frac c3 = 0$ कहता है; $a$ के लिए हल कीजिए और
प्राचलन कीजिए:
\[
P = b\Bigl(X - \frac12\Bigr) + c\Bigl(X^2 - \frac13\Bigr),
\]
अतः $H = \operatorname{Vect}\bigl(X - \frac12,\ X^2 - \frac13\bigr)$, विमा
$2$ (दोनों बहुपदों की घातें भिन्न हैं: \omterm{def:b1:vspaces:free}{स्वतंत्र})। कोई एक रैखिक शर्त --- चाहे
वह मूल्यांकन हो, \omterm{thm:b1:integration:def}{समाकल} हो, या कोई और रैखिक नुस्ख़ा --- अधिक से अधिक एक विमा
हटाती है, और शर्त के सर्वसम शून्य न होते ही ठीक एक। \cref{ch:b1:linmaps} ऐसे
नुस्ख़ों को \emph{रैखिक रूप} और उनके हल-समुच्चयों को \emph{अधिसमतल} नाम
देगा; और यहाँ मिले आधार-सदिश \cref{ch:b1:euclid} में लजांद्र कुल के आरंभ के
रूप में फिर से दिखाई देंगे।
\end{example}

\section{उपसमष्टियाँ, कोटि, ग्रासमान}

\begin{theorem}[उपसमष्टियाँ]\label{thm:b1:findim:subspaces}
मान लीजिए $E$ \omterm{def:b1:findim:def}{परिमित-विमीय} है और $F$ उसकी \omterm{def:b1:vspaces:subspace}{उपसमष्टि}। तब $F$ \omterm{def:b1:findim:def}{परिमित-विमीय} है,
$\dim F \leq \dim E$, और समानता तभी होती है जब $F = E$। इसके अतिरिक्त हर
\omterm{def:b1:vspaces:subspace}{उपसमष्टि} की कोई \omterm{def:b1:vspaces:sum}{पूरक} \omterm{def:b1:vspaces:subspace}{उपसमष्टि} होती है।
\end{theorem}

\begin{proof}
$F$ के \omterm{def:b1:vspaces:free}{स्वतंत्र} कुलों में अधिक से अधिक $\dim E$ सदिश होते हैं ($E$ में
विनिमय प्रमेयिका: $F$ का कोई \omterm{def:b1:vspaces:free}{स्वतंत्र कुल} विशेष रूप से $E$ में भी \omterm{def:b1:vspaces:free}{स्वतंत्र}
है, और $E$ का कोई परिमित जनक कुल है)। $F$ के \omterm{def:b1:vspaces:free}{स्वतंत्र} कुलों में से
\emph{अधिकतम} आकार $p$ वाला एक चुनिए --- यह संभव है, क्योंकि आकार $\dim E$
से परिबद्ध पूर्णांक हैं। वह $F$ को जनित करता है: अन्यथा कोई $x \in F$ उसकी
\omterm{def:b1:vspaces:span}{जनित समष्टि} के बाहर होता, और $x$ जोड़ने पर $F$ का $p + 1$ आकार वाला \omterm{def:b1:vspaces:free}{स्वतंत्र
कुल} मिल जाता (\cref{prop:b1:vspaces:freecriteria}~(2)), जो अधिकतमता का विरोध
करता। \omterm{def:b1:vspaces:free}{स्वतंत्र} और जनक होने के कारण वह $F$ का \omterm{def:b1:vspaces:free}{आधार} है, और $\dim F = p \leq
\dim E$। यदि $p = \dim E = n$: तो $E$ के $n$ सदिशों वाला \omterm{def:b1:vspaces:free}{स्वतंत्र कुल} $E$ का
\omterm{def:b1:vspaces:free}{आधार} है (\cref{prop:b1:findim:twoofthree}), अतः $F \supseteq
\operatorname{span} = E$। \omterm{def:b1:vspaces:sum}{पूरक}: $F$ के \omterm{def:b1:vspaces:free}{आधार} $(f_1, \dots, f_p)$ को अपूर्ण
\omterm{def:b1:vspaces:free}{आधार} प्रमेय से $E$ के \omterm{def:b1:vspaces:free}{आधार} $(f_1, \dots, f_p, g_{p+1}, \dots, g_n)$ तक पूर्ण
कीजिए; तब $G = \operatorname{Vect}(g_{p+1}, \dots, g_n)$ $E = F \oplus G$
संतुष्ट करती है (विघटनों का अस्तित्व और अद्वितीयता = बड़े \omterm{def:b1:vspaces:free}{आधार} में
\omterm{prop:b1:vspaces:coordinates}{निर्देशांक})।
\end{proof}

\begin{definition}[किसी कुल की कोटि]\label{def:b1:findim:rank}
सदिशों के किसी परिमित कुल की \emph{कोटि}\index{कोटि!किसी कुल की} उसकी \omterm{def:b1:vspaces:span}{जनित
समष्टि} की विमा है: $\operatorname{rk}(x_1, \dots, x_p) = \dim
\operatorname{Vect}(x_1, \dots, x_p) \leq \min(p, \dim E)$, और $p$ के साथ
समानता तभी होती है जब कुल \omterm{def:b1:vspaces:free}{स्वतंत्र} हो।
\end{definition}

\begin{example}[विलोपन से कोटि की संगणना]
\label{ex:b1:findim:rankcomputation}
$\R^3$ में $\bigl((1,2,3), (2,3,4), (3,4,5), (1,1,1)\bigr)$ की \omterm{def:b1:findim:rank}{कोटि}। किसी
सदिश में से शेष का संयोजन घटाने पर \omterm{def:b1:vspaces:span}{जनित समष्टि} नहीं बदलती (दोनों कुल एक ही
संयोजन जनित करते हैं): $(2,3,4)$ के स्थान पर $(2,3,4) - (1,2,3) = (1,1,1)$
और $(3,4,5)$ के स्थान पर $(3,4,5) - (1,2,3) = (2,2,2)$ रखिए। अब \omterm{def:b1:vspaces:span}{जनित समष्टि}
$\operatorname{Vect}\bigl((1,2,3), (1,1,1), (2,2,2), (1,1,1)\bigr) =
\operatorname{Vect}\bigl((1,2,3), (1,1,1)\bigr)$ है, और ये दोनों सदिश
अनुपाती नहीं हैं: अतः \omterm{def:b1:findim:rank}{कोटि} $2$ है। घटाओ-और-हटाओ वाली यह प्रक्रिया
\cref{ch:b1:det} में गाउसीय विलोपन के रूप में व्यवस्थित की गई है।
\end{example}

\begin{example}[जनकों को एक साथ लिखकर योग]
\label{ex:b1:findim:sumconcat}
$\R^3$ में लीजिए
\[
F = \operatorname{Vect}\bigl((1,2,3),\ (1,1,1)\bigr),
\qquad
G = \operatorname{Vect}\bigl((2,3,4)\bigr) .
\]
योग $F + G$ तीनों जनकों को एक साथ लिखने से बने कुल द्वारा जनित है, और
\[
(2, 3, 4) = (1, 2, 3) + (1, 1, 1)
\]
यह दिखाता है कि तीसरा अनावश्यक है: $F + G = F$, जिसकी विमा $2$ है --- तुल्य
रूप से $G \subseteq F$, जिसे यह संबंध दिखा देता है। ग्रासमान पुष्टि करता है:
$\dim(F \cap G) = 2 + 1 - 2 = 1 = \dim G$। योग जनकों को एक साथ लिखकर और फिर
उस ढेर को कोटि-कलनविधि से घटाकर संगणित किए जाते हैं; कोई नई तकनीक कभी नहीं
चाहिए।
\end{example}

\begin{theorem}[ग्रासमान का सूत्र]\label{thm:b1:findim:grassmann}
$E$ की \omterm{def:b1:findim:def}{परिमित-विमीय} उपसमष्टियों $F, G$ के लिए:\index{ग्रासमान का सूत्र}
\[
\dim(F + G) = \dim F + \dim G - \dim(F \cap G) .
\]
विशेष रूप से $F + G$ प्रत्यक्ष है तभी जब $\dim(F + G) = \dim F + \dim G$।
\end{theorem}

\begin{proof}
$F \cap G$ के \omterm{def:b1:vspaces:free}{आधार} $(e_1, \dots, e_r)$ से आरंभ कीजिए; उसे $F$ के \omterm{def:b1:vspaces:free}{आधार} $(e_1,
\dots, e_r, f_1, \dots, f_s)$ तक और $G$ के \omterm{def:b1:vspaces:free}{आधार} $(e_1, \dots, e_r, g_1,
\dots, g_t)$ तक पूर्ण कीजिए (\cref{thm:b1:findim:bases}~(2))। हमारा दावा है
कि
\[
\mathcal{B} = (e_1, \dots, e_r, f_1, \dots, f_s, g_1, \dots, g_t)
\]
$F + G$ का \omterm{def:b1:vspaces:free}{आधार} है; और सूत्र गिनती से आ जाता है: $(r + s) + (r + t) - r = r
+ s + t$।

$\mathcal{B}$ $F + G$ को जनित करता है: कोई भी $u + v$ ($u \in F$, $v \in G$)
उसके द्वारा प्रसारित हो जाता है। स्वतंत्रता: मान लीजिए $\sum \alpha_i e_i +
\sum \beta_j f_j + \sum \gamma_k g_k = 0$। सदिश $w = \sum \gamma_k g_k =
-\sum\alpha_i e_i - \sum\beta_j f_j$ $G \cap F$ में है, अतः वह अकेले $(e_i)$
पर प्रसारित होता है; पर $w$ $(g_k)$ पर भी प्रसारित होता है, और $G$ के \omterm{def:b1:vspaces:free}{आधार}
में ये दोनों व्यंजक एक ही होने चाहिए: अतः सभी $\gamma_k = 0$ (और
$e$-निर्देशांक मेल खाते हैं)। तब संबंध घटकर $\sum\alpha_i e_i + \sum\beta_j
f_j = 0$ रह जाता है, जो $F$ के \omterm{def:b1:vspaces:free}{आधार} में एक संबंध है: शेष सभी गुणांक शून्य हो
जाते हैं।
\end{proof}

\begin{example}\label{ex:b1:findim:planes}
$\R^3$ के दो भिन्न समतल $F, G$ (विमा $2$) $F + G = \R^3$ संतुष्ट करते हैं
(उनका योग किसी समतल को सचमुच समाहित करता है), अतः $\dim(F \cap G) = 2 + 2 -
3 = 1$: वे सदा किसी रेखा के अनुदिश काटते हैं --- मूल बिंदु से होकर जाने वाले
``समांतर समतल'' होते ही नहीं।
\end{example}

\begin{example}[ग्रासमान काम पर, $\R^4$ में]
\label{ex:b1:findim:grassmannrun}
$\R^4$ में मान लीजिए $F = \operatorname{Vect}\bigl(e_1,\ e_2,\
(1,1,1,0)\bigr)$ और $G = \operatorname{Vect}(e_3, e_4)$। विमाएँ: $\dim F =
3$ (तीसरे जनक का तीसरा \omterm{prop:b1:vspaces:coordinates}{निर्देशांक} अशून्य है, अतः वह
$\operatorname{Vect}(e_1, e_2)$ के बाहर है) और $\dim G = 2$। योग: $F + G$
में $e_1, e_2, e_4$ और $e_3 = (1,1,1,0) - e_1 - e_2$ हैं: अतः वह पूरा $\R^4$
है। तब ग्रासमान बिना किसी विलोपन के सर्वनिष्ठ का आकार \emph{संगणित} कर देता
है:
\[
\dim(F \cap G) = 3 + 2 - 4 = 1 .
\]
रेखा पहचानने के लिए $G$ के भीतर देखिए: सदिश $(0, 0, c, d)$ $F$ में ठीक तब है
जब वह $a e_1 + b e_2 + \lambda(1,1,1,0)$ हो, जो $\lambda = -a = -b$ और $d =
0$ बाध्य कर देता है: अतः सर्वनिष्ठ $\operatorname{Vect}(e_3)$ है --- और यह
संगत है, क्योंकि $e_3$ ऊपर $F$ में दिखाया गया था और परिभाषा से $G$ में है।
श्रम का प्रतिरूपी \omterm{thm:b1:logic:partition}{विभाजन}: ग्रासमान बताता है कि \emph{कितना} खोजना है, और फिर
रैखिक निकाय खोज लेता है कि \emph{क्या}।
\end{example}

\begin{example}[समीकरणों से सर्वनिष्ठ निकालना]
\label{ex:b1:findim:intersectequations}
जब दोनों उपसमष्टियाँ \omterm{def:b1:logic:sets}{हल-समुच्चय} के रूप में आती हों, तो सर्वनिष्ठ निकालना बस
समीकरणों को एक के ऊपर एक रख देना है। $\R^3$ में: $F = \{x + y + z = 0\}$ और
$G = \{x = y\}$ से मिलता है
\[
F \cap G = \{x = y,\ 2x + z = 0\}
= \{(x,\ x,\ -2x)\} = \operatorname{Vect}(1, 1, -2),
\]
अर्थात् एक रेखा। ग्रासमान से जाँच: $F + G = \R^3$ (समतल भिन्न हैं, अतः उनका
योग किसी समतल को सचमुच समाहित करता है), अतः $\dim(F\cap G) = 2 + 2 - 3 = 1$।
किसी \omterm{def:b1:vspaces:subspace}{उपसमष्टि} के दोनों वर्णन --- समीकरणों द्वारा और जनकों द्वारा --- में से
हर एक किसी एक संक्रिया को तुच्छ बना देता है: समीकरण एक के ऊपर एक रखकर काटते
हैं, जनक एक साथ लिखकर जुड़ते हैं; और इन दोनों के बीच अनुवाद करना ही रैखिक
निकाय हल करने का अर्थ है (\cref{ch:b1:det})।
\end{example}

\begin{remark}[सामान्य भूलें]
\label{rem:b1:findim:pitfalls}
\emph{योगों के अनुदिश विमाएँ जुड़ती नहीं} जब तक योग प्रत्यक्ष न हो: $\R^3$
के दो समतलों के लिए $\dim(F + G) = 3$ होता है, $4$ नहीं; सदा सर्वनिष्ठ वाले
पद से सुधार कीजिए (ग्रासमान)। \emph{अंतर्विष्टता के लिए विमा \emph{और}
अंतर्विष्टता, दोनों चाहिए}: अकेला $\dim F = \dim G$ कभी $F = G$ नहीं देता
($\R^2$ की दो भिन्न रेखाएँ); \cref{thm:b1:findim:subspaces} की समानता वाली
स्थिति के लिए पहले $F \subseteq G$ चाहिए। \emph{प्राचल गिन लेना अभी उपपत्ति
नहीं है}: ``$\R^4$ में दो समीकरण, अतः विमा $2$'' तब विफल हो जाता है जब
समीकरण परतंत्र हों ($x + y = 0$ और $2x + 2y = 0$ विमा $3$ छोड़ते हैं);
निर्णय केवल कोई वास्तविक प्राचलन या \omterm{def:b1:findim:rank}{कोटि} की संगणना ही करती है। \emph{जो
\omterm{def:b1:vspaces:subspace}{उपसमष्टि} न हो, उसकी विमा की बात मत कीजिए}: \emph{असमघातीय} निकायों के
हल-समुच्चयों में $0$ नहीं होता; उनकी ``विमा'' संबद्ध समघातीय हल-समष्टि की
विमा है (\cref{ch:b1:det} इसे यथार्थ बना देता है)। \emph{अनंत विमा होती है}:
$K[X]$ में हर आकार के \omterm{def:b1:vspaces:free}{स्वतंत्र कुल} हैं (एकपदी), अतः कोई परिमित जनक कुल हो ही
नहीं सकता --- तीन में से दो वाले नियम जैसे \omterm{def:b1:logic:statement}{कथन} कड़ाई से \omterm{def:b1:findim:def}{परिमित-विमीय} हैं और
$K[X]$ पर बुरी तरह विफल हो जाते हैं (\cref{cor:b1:linmaps:samedim} \omterm{def:b1:logic:inj}{एकैकी} तथा
\omterm{def:b1:logic:inj}{आच्छादक} के लिए यही दिखाएगा)। \emph{\omterm{def:b1:findim:rank}{कोटि} \omterm{def:b1:vspaces:span}{जनित समष्टि} की बात है, सूची की
नहीं}: किसी सदिश को दोहराने, क्रम बदलने या अशून्य अचरों से मापित करने पर
\omterm{def:b1:findim:rank}{कोटि} नहीं बदलती, और $\operatorname{rk} = p$ (सदिशों की संख्या) एक
\emph{सिद्ध करने योग्य गुण} है --- वह ठीक स्वतंत्रता ही है। $2$ \omterm{def:b1:findim:rank}{कोटि} वाले
$5$ सदिशों के कुल में तीन सदिशों जितनी अतिरेकता है, जिसे विलोपन
(\cref{ex:b1:findim:rankcomputation}) स्पष्ट रूप से खोज देता है।
\end{remark}

\begin{remark}[विमा कहाँ काम पर जाती है]
\label{rem:b1:findim:whereused}
यहाँ से आगे विमा इस पुस्तक का सबसे प्रिय गिनती-तर्क है। \cref{ch:b1:linmaps}
कोटि--शून्यता प्रमेय सिद्ध करता है, जो ग्रासमान के सूत्र का फलनात्मक रूप है;
\cref{ch:b1:matrices} पंक्ति-लघुकरण से कोटियाँ संगणित करता है;
\cref{ch:b1:det} ``$K^n$ के $n$ सदिश \omterm{def:b1:vspaces:free}{आधार} बनाते हैं'' को एक ही संख्या के
अशून्य होने में बदल देता है। नीचे की सप्ताहांत समस्या विमा को \emph{अंकगणित}
करते हुए दिखाती है: \omterm{def:b1:structures:field}{क्षेत्र} $\Q$ पर विमाएँ गिनकर ऐसे अपरिमेयता-कथन सिद्ध हो
जाते हैं जो हाथ से अछूते लगते हैं। स्नातक वर्ष~3 के खंड में वही
विमा-गिनतियाँ, \omterm{def:b1:structures:group}{समूह} सिद्धांत से परिष्कृत होकर, तय करती हैं कि कौन-सी
शास्त्रीय रचना-समस्याएँ हल हो सकती हैं --- वह कथा गैलुआ सिद्धांत है।
\end{remark}

\begin{remark}[पुस्तक 3 के भीतर के परिदृश्य]
\label{rem:b1:findim:perspectives}
इस खंड के शेष भाग में विमा ही संरक्षित राशि है, और संरक्षण के नियमों को पहले
ही नाम दे देना उचित है। \cref{ch:b1:linmaps} $\dim E = \dim\ker u +
\operatorname{rk} u$ सिद्ध करता है: कोई रैखिक \omterm{def:b1:logic:map}{प्रतिचित्रण} जो कुचल देता है और
जो बचा रखता है, दोनों का योग सदा स्रोत के बराबर होता है। \cref{ch:b1:det}
इसे हल-समुच्चयों की संरचना तक परिष्कृत करता है: $p$ अज्ञातों में से
$\operatorname{rk} A$ धुरियाँ घटाने पर हल-समष्टि की विमा बचती है, जिसे
गाउसीय विलोपन मुक्त प्राचलों के रूप में दिखा देता है। \cref{ch:b1:euclid}
$\dim E = \dim F + \dim F^\perp$ का लांबिक \omterm{thm:b1:logic:partition}{विभाजन} करता है, और
\cref{ch:b1:multivar} की सप्ताहांत समस्या ठीक यही बजट ख़र्च करती है: $n$
आँकड़ा-बिंदु, $2$ आरोपित प्राचल, $n - 2$ विमाएँ अवशिष्ट की। आगे किसी अध्याय
में जब भी कोई गिनती संतुलित होने से इनकार करे, तो त्रुटि कोई भूली हुई अष्टि
है या कोई अप्रत्यक्ष योग --- सबसे पहले ग्रासमान के सूत्र पर लौटिए।
\end{remark}

\section{अभ्यास}

\begin{exercise}[$\star$]\label{exo:b1:findim:1}
निम्नलिखित का एक \omterm{def:b1:vspaces:free}{आधार} और विमा दीजिए:
\begin{enumerate}
  \item $F = \{(x,y,z) \in \R^3 : x + y + z = 0\}$;
  \item $G = \{(x,y,z,t) \in \R^4 : x = y,\ z = 2t\}$;
  \item $H = \{P \in \R_3[X] : P(1) = P'(1) = 0\}$.
\end{enumerate}
\end{exercise}

\begin{exercise}[$\star$]\label{exo:b1:findim:2}
$\R^3$ में कुल $\bigl((1,1,1), (1,2,3), (3,5,7), (0,1,2)\bigr)$ की \omterm{def:b1:findim:rank}{कोटि}
संगणित कीजिए, और उसकी \omterm{def:b1:vspaces:span}{जनित समष्टि} का एक \omterm{def:b1:vspaces:free}{आधार} निकालिए।
\end{exercise}

\begin{exercise}[$\star$]\label{exo:b1:findim:3}
विहित सदिशों का प्रयोग करके \omterm{def:b1:vspaces:free}{स्वतंत्र कुल} $\bigl((1,1,0,0), (0,0,1,1)\bigr)$
को $\R^4$ के \omterm{def:b1:vspaces:free}{आधार} तक पूर्ण कीजिए, और औचित्य दीजिए।
\end{exercise}

\begin{exercise}[$\star$]\label{exo:b1:findim:4}
सिद्ध कीजिए कि $\bigl(1, X, X(X-1), X(X-1)(X-2)\bigr)$ $\R_3[X]$ का \omterm{def:b1:vspaces:free}{आधार} है,
और उसमें $X^3$ के \omterm{prop:b1:vspaces:coordinates}{निर्देशांक} खोजिए।
\end{exercise}

\begin{exercise}[$\star\star$]\label{exo:b1:findim:5}
मान लीजिए $F$ और $G$ किसी समष्टि $E$ की क्रमशः $4$ और $5$ विमा वाली
उपसमष्टियाँ हैं, जहाँ $\dim E = 7$। $\dim(F \cap G)$ के संभव मान क्या हैं?
$E = \R^7$ के साथ हर मान को साकार करने वाला एक उदाहरण दीजिए।
\end{exercise}

\begin{exercise}[$\star\star$]\label{exo:b1:findim:6}
मान लीजिए $H = \{P \in \R_n[X] : P(1) = 0\}$। सिद्ध कीजिए कि $H$ $\R_n[X]$
का अधिसमतल है ($n$ विमा वाली \omterm{def:b1:vspaces:subspace}{उपसमष्टि}), $H$ का एक \omterm{def:b1:vspaces:free}{आधार} दिखाइए \emph{(गुणनखंड
प्रमेय याद कीजिए: $P = (X-1)Q$)}, और एक \omterm{def:b1:vspaces:sum}{पूरक} रेखा दीजिए।
\end{exercise}

\begin{exercise}[$\star\star$]\label{exo:b1:findim:7}
मान लीजिए $u_1, \dots, u_p$ $r$ \omterm{def:b1:findim:rank}{कोटि} वाले सदिश हैं। सिद्ध कीजिए कि एक सदिश
हटाने पर $r$ या $r - 1$ \omterm{def:b1:findim:rank}{कोटि} वाला कुल मिलता है, और एक सदिश जोड़ने पर \omterm{def:b1:findim:rank}{कोटि}
$r$ या $r + 1$ होती है। इससे निष्कर्ष निकालिए कि किसी भी एक निवेश या विलोपन
से \omterm{def:b1:findim:rank}{कोटि} अधिक से अधिक $1$ बदलती है।
\end{exercise}

\begin{exercise}[$\star\star\star$]\label{exo:b1:findim:8}
मान लीजिए $F_1 \subseteq F_2 \subseteq \dots \subseteq F_k$ $E$ ($\dim E =
n$) की ऐसी उपसमष्टियाँ हैं कि सभी $i$ के लिए $F_i \neq F_{i+1}$। सिद्ध कीजिए
$k \leq n + 1$। इससे निष्कर्ष निकालिए कि $\R^n$ की उपसमष्टियों की कोई भी
सचमुच वर्धमान शृंखला अधिक से अधिक $n + 1$ लंबी होती है, और अधिकतम लंबाई वाली
एक शृंखला दिखाइए।
\end{exercise}

\begin{exercise}[$\star\star\star$]\label{exo:b1:findim:9}
मान लीजिए $E$ की विमा $n$ है और $F$, $G$ दो अधिसमतल हैं (विमा $n - 1$), $F
\neq G$। $\dim(F \cap G)$ संगणित कीजिए। सामान्यीकरण कीजिए: $k$ अधिसमतलों के
सर्वनिष्ठ की विमा $\geq n - k$ होती है।
\end{exercise}

\begin{exercise}[$\star\star$]\label{exo:b1:findim:10}
मान लीजिए $E$ उन वास्तविक अनुक्रमों का \omterm{def:b1:logic:sets}{समुच्चय} है जो सभी $n$ के लिए $u_{n+2}
= 3u_{n+1} - 2u_n$ संतुष्ट करते हैं।
\begin{enumerate}
  \item दिखाइए कि $E$ अनुक्रमों की समष्टि की \omterm{def:b1:vspaces:subspace}{उपसमष्टि} है, और यह कि $E$ का
  कोई अनुक्रम युग्म $(u_0, u_1)$ से रैखिक रूप से पूर्णतः तय हो जाता है; इससे
  $\dim E = 2$ निकालिए।
  \item जाँचिए कि अचर अनुक्रम $(1)$ और गुणोत्तर अनुक्रम $(2^n)$ $E$ में हैं
  और $E$ का \omterm{def:b1:vspaces:free}{आधार} बनाते हैं।
  \item $u_0 = 0$, $u_1 = 1$ वाला $E$ का अनुक्रम खोजिए।
\end{enumerate}
\end{exercise}

\begin{exercise}[$\star\star$]\label{exo:b1:findim:11}
मान लीजिए $E$ की विमा $n$ है।
\begin{enumerate}
  \item यदि $F, G$ ऐसी उपसमष्टियाँ हैं कि $\dim F + \dim G > n$, तो सिद्ध
  कीजिए $F \cap G \neq \{0\}$। उदाहरण दीजिए: $\R^{100}$ की $51$ और $50$ विमा
  वाली दो उपसमष्टियों में सदा कोई अशून्य सदिश उभयनिष्ठ होता है।
  \item यदि $H$ कोई अधिसमतल है और $F$ ऐसी \omterm{def:b1:vspaces:subspace}{उपसमष्टि} कि $F \cap H = \{0\}$, तो
  सिद्ध कीजिए $\dim F \leq 1$।
\end{enumerate}
\end{exercise}

\begin{exercise}[$\star\star\star$]\label{exo:b1:findim:12}
(उभयनिष्ठ \omterm{def:b1:vspaces:sum}{पूरक}) मान लीजिए $F, G$ $E$ (\omterm{def:b1:findim:def}{परिमित विमा}) की ऐसी उपसमष्टियाँ हैं कि
$\dim F = \dim G$। सिद्ध कीजिए कि $F$ और $G$ का कोई \emph{उभयनिष्ठ} \omterm{def:b1:vspaces:sum}{पूरक} है:
अर्थात् ऐसी \omterm{def:b1:vspaces:subspace}{उपसमष्टि} $S$ है कि $E = F \oplus S = G \oplus S$। \emph{($\dim
F$ पर नीचे की ओर आगमन कीजिए: यदि $F \neq G \neq E$, तो $x \notin F \cup G$
चुनिए --- \cref{exo:b1:vspaces:12} इसकी अनुमति देता है --- और $F \oplus
\operatorname{Vect}(x)$ तथा $G \oplus \operatorname{Vect}(x)$ पर विचार
कीजिए।)}
\end{exercise}

\section{समस्या: डेडेकिंड का मीनार नियम}

\begin{problem}\label{pb:b1:findim:1}
अध्याय 18--19 में अदिशों के विषय में क्षेत्र-अभिगृहीतों
(\cref{def:b1:structures:field}) से आगे कुछ भी प्रयुक्त नहीं हुआ: अतः $K =
\Q$ लिया जा सकता है और \emph{वास्तविक संख्याओं} के \omterm{def:b1:logic:sets}{समुच्चय} $\Q$-विमा के गज़
से नापे जा सकते हैं। यह समस्या $\Q(\sqrt2, \sqrt3)$ की विमा संगणित करती है,
\emph{डेडेकिंड का \omterm{pb:b1:findim:1}{मीनार नियम}} $\dim_\Q M = \dim_\Q K \cdot \dim_K M$ सिद्ध
करती है, और शुद्ध विमा-गिनती से अपरिमेयता की प्रमेयें काट लेती है --- न
$\varepsilon$, न दशमलव, केवल \omterm{def:b1:vspaces:free}{आधार}।\index{मीनार नियम} \index{क्षेत्र!$\R$ का
उपक्षेत्र}\index{अपरिमेय संख्या}

\medskip \noindent\textbf{भाग I --- परिमेय अदिश।}
\begin{enumerate}
  \item जाँचिए कि $\Q$ एक \omterm{def:b1:structures:field}{क्षेत्र} है और अध्याय 18--19 की हर परिभाषा तथा
  उपपत्ति अदिशों के केवल क्षेत्र-अभिगृहीतों का प्रयोग करती है; निष्कर्ष
  निकालिए कि $\R$ एक $\Q$-सदिश समष्टि है और विनिमय प्रमेयिका, आधार-प्रमेयें
  तथा ग्रासमान का सूत्र $\Q$ पर लागू होते हैं। \cref{thm:b1:findim:exchange}
  की उपपत्ति का वह एक चरण बताइए जहाँ किसी अशून्य अदिश से भाग दिया जाता है।
  \item दिखाइए कि $(1, \sqrt2)$ $\Q$ पर \omterm{def:b1:vspaces:free}{स्वतंत्र} है पर $\R$ पर परतंत्र।
  $\Q(\sqrt2) = \operatorname{Vect}_\Q(1, \sqrt2) = \{a + b\sqrt2 : a, b \in
  \Q\}$ रखिए; $\dim_\Q \Q(\sqrt2)$ क्या है?
  \item दिखाइए कि $\Q(\sqrt2)$ गुणन के अंतर्गत स्थायी है। $u = a + b\sqrt2$
  के लिए $\sigma(u) = a - b\sqrt2$ और $N(u) = u\,\sigma(u) = a^2 - 2b^2$
  रखिए। $\sigma(uv) = \sigma(u)\sigma(v)$ दिखाइए, $N(uv) = N(u)N(v)$
  निकालिए, और दिखाइए कि $u \neq 0$ होने पर $N(u) \neq 0$।
  \item इससे निष्कर्ष निकालिए कि हर अशून्य $u \in \Q(\sqrt2)$ का प्रतिलोम
  $\Q(\sqrt2)$ में ही है, अर्थात् $u^{-1} = \sigma(u)/N(u)$: अतः
  $\Q(\sqrt2)$ $\R$ का उपक्षेत्र है। $\dfrac1{3 + 2\sqrt2}$ और $\dfrac1{1 +
  \sqrt2}$ संगणित कीजिए।
\end{enumerate}

\medskip \noindent\textbf{भाग II --- $\sqrt3$ को जोड़ना।}
\begin{enumerate}[resume]
  \item सिद्ध कीजिए कि $\sqrt6 \notin \Q$ \emph{($6q^2 = p^2$ के दोनों
  पक्षों पर $2$ के घातांक की तुलना कीजिए, जैसा \cref{exo:b1:arith:7} में)},
  फिर यह कि $\sqrt3 \notin \Q(\sqrt2)$ \emph{($\sqrt3 = a + b\sqrt2$ का वर्ग
  कीजिए और $ab \neq 0$, $b = 0$, $a = 0$ स्थितियों पर चर्चा कीजिए)}।
  \item मान लीजिए $M = \operatorname{Vect}_\Q(1, \sqrt2, \sqrt3, \sqrt6)$।
  दिखाइए कि $M$ गुणन के अंतर्गत स्थायी है \emph{(आधार-सदिशों के गुणनफलों की
  एक सारणी पर्याप्त है)}।
  \item $K = \Q(\sqrt2)$ लिखिए। दिखाइए कि $(1, \sqrt3)$ \emph{\omterm{def:b1:structures:field}{क्षेत्र} $K$
  पर} \omterm{def:b1:vspaces:free}{स्वतंत्र} है, और निष्कर्ष निकालिए कि $M = K + K\sqrt3$ विमा $2$ वाली
  $K$-सदिश समष्टि है जिसका \omterm{def:b1:vspaces:free}{आधार} $(1, \sqrt3)$ है।
  \item सिद्ध कीजिए कि $(1, \sqrt2, \sqrt3, \sqrt6)$ $\Q$ पर \omterm{def:b1:vspaces:free}{स्वतंत्र} है,
  अतः $\dim_\Q M = 4$। \emph{(किसी शून्य संबंध को $(a + b\sqrt2) + (c +
  d\sqrt2)\sqrt3 = 0$ के रूप में समूहबद्ध कीजिए और पहले प्रश्न 7, फिर प्रश्न
  2 लगाइए।)}
\end{enumerate}

\medskip \noindent\textbf{भाग III --- \omterm{pb:b1:findim:1}{मीनार नियम}।} मान लीजिए $\Q \subseteq K
\subseteq M \subseteq \R$, जहाँ $K$ और $M$ उपक्षेत्र हैं, $(e_1, \dots,
e_m)$ $\Q$-सदिश समष्टि के रूप में $K$ का \omterm{def:b1:vspaces:free}{आधार} है, और $(f_1, \dots, f_n)$
$K$-सदिश समष्टि के रूप में $M$ का \omterm{def:b1:vspaces:free}{आधार}।
\begin{enumerate}[resume]
  \item दिखाइए कि $mn$ गुणनफल $(e_i f_j)$ $\Q$ पर $M$ को जनित करते हैं।
  \item दिखाइए कि कुल $(e_i f_j)$ $\Q$ पर \omterm{def:b1:vspaces:free}{स्वतंत्र} है। \emph{(किसी शून्य
  $\Q$-संयोजन को $\sum_j \bigl(\sum_i \lambda_{ij} e_i\bigr) f_j$ के रूप में
  पुनर्व्यवस्थित कीजिए, जिसके भीतरी गुणांक $K$ में रहते हैं।)}
  \item \emph{डेडेकिंड के \omterm{pb:b1:findim:1}{मीनार नियम}} के साथ निष्कर्ष निकालिए:
        \[
        \dim_\Q M \;=\; \dim_\Q K \,\cdot\, \dim_K M ,
        \]
        और उसे $\Q \subseteq \Q(\sqrt2) \subseteq M$ पर प्रश्न 7 तथा 8 के
        सामने जाँचिए।
  \item (मुफ़्त में \omterm{def:b1:structures:field}{क्षेत्र}) मान लीजिए $A \subseteq \R$ $1$ को समाहित करने
  वाली, गुणन के अंतर्गत स्थायी, \omterm{def:b1:findim:def}{परिमित-विमीय} $\Q$-उपसमष्टि है, और $u \in A$,
  $u \neq 0$। दिखाइए कि यदि $(a_1, \dots, a_d)$ $A$ का \omterm{def:b1:vspaces:free}{आधार} है, तो $(u a_1,
  \dots, u a_d)$ भी $A$ का \omterm{def:b1:vspaces:free}{आधार} है; इससे निष्कर्ष निकालिए कि $u$ का प्रतिलोम
  \emph{$A$ में ही} है: अतः $A$ $\R$ का उपक्षेत्र है। इससे पहले के कौन-से
  प्रश्न फिर से मिल जाते हैं?
  \item दिखाइए कि हर $x \in A$ के लिए (प्रश्न 12 की तरह, $\dim_\Q A = d$)
  कुल $(1, x, x^2, \dots, x^{d})$ परतंत्र है: अर्थात् $A$ का हर अवयव परिमेय
  गुणांकों वाले किसी अशून्य \omterm{def:b1:poly:def}{बहुपद} का मूल है, जिसकी घात अधिक से अधिक $d$ है।
\end{enumerate}

\medskip \noindent\textbf{भाग IV --- एक ही संख्या सब कुछ जनित कर देती है।}
$s = \sqrt2 + \sqrt3$ रखिए।
\begin{enumerate}[resume]
  \item $M$ के \omterm{def:b1:vspaces:free}{आधार} $(1, \sqrt2, \sqrt3, \sqrt6)$ में $s^2$, $s^3$ और $s^4$
  के \omterm{prop:b1:vspaces:coordinates}{निर्देशांक} संगणित कीजिए।
  \item दिखाइए कि $(1, s, s^2, s^3)$ $\Q$ पर \omterm{def:b1:vspaces:free}{स्वतंत्र} है, और
  $\operatorname{Vect}_\Q(1, s, s^2, s^3) = M$ निकालिए: $M$ का हर अवयव $s$
  में एक परिमेय \omterm{def:b1:poly:def}{बहुपद} है। $\sqrt2$ और $\sqrt3$ को ऐसे ही बहुपदों के रूप में
  व्यक्त कीजिए।
  \item $s^4 - 10s^2 + 1 = 0$ सत्यापित कीजिए, और दिखाइए कि $X^4 - 10X^2 + 1$
  $s$ पर शून्य होने वाला \emph{न्यूनतम घात का \omterm{def:b1:poly:def}{इकाई-अग्र बहुपद}} है। उसके
  चारों वास्तविक मूल निर्धारित कीजिए।
  \item निष्कर्ष निकालिए: $1/s = 10s - s^3$; इस संख्या को पहचानिए। दिखाइए कि
  $a + b\sqrt2 + c\sqrt3 + d\sqrt6$ ($a,b,c,d \in \Q$) परिमेय है तभी जब $b =
  c = d = 0$; विशेष रूप से $\sqrt2 + \sqrt3 + \sqrt6$ अपरिमेय है
  (\cref{exo:b1:reals:7} को और सुदृढ़ करते हुए)।
\end{enumerate}

\medskip \noindent\textbf{भाग V --- घनमूल बाहर ही रहता है।} $t = 2^{1/3}$
रखिए।
\begin{enumerate}[resume]
  \item दिखाइए कि $X^3 - 2$ का कोई परिमेय मूल नहीं है
  \emph{(\cref{exo:b1:poly:5} की परिमेय-मूल कसौटी, अथवा कोई
  मूल्यांकन-गिनती)}; इससे निष्कर्ष निकालिए कि $(1, t)$ $\Q$ पर \omterm{def:b1:vspaces:free}{स्वतंत्र} है।
  \item दिखाइए कि $(1, t, t^2)$ $\Q$ पर \omterm{def:b1:vspaces:free}{स्वतंत्र} है। \emph{(यदि $\leq 2$ घात
  वाला कोई अशून्य $P \in \Q[X]$ $t$ को मार देता है, तो न्यूनतम घात वाला एक
  लीजिए और $X^3 - 2$ को उससे भाग दीजिए, \cref{thm:b1:poly:division};
  निष्कर्ष निकालिए कि $X^3 - 2$ का कोई परिमेय मूल होता।)} इससे निकालिए कि $A
  = \operatorname{Vect}_\Q(1, t, t^2)$ की विमा $3$ है, वह गुणन के अंतर्गत
  स्थायी है, और $\R$ का उपक्षेत्र है।
  \item सिद्ध कीजिए कि $t \notin M$: अन्यथा $M$ \omterm{def:b1:structures:field}{क्षेत्र} $A$ पर एक \omterm{def:b1:vspaces:def}{सदिश
  समष्टि} होती, और \omterm{pb:b1:findim:1}{मीनार नियम} $3 \mid 4$ बाध्य कर देता। अतः $2^{1/3}$ $1,
  \sqrt2, \sqrt3, \sqrt6$ का परिमेय संयोजन नहीं है।
  \item इससे निष्कर्ष निकालिए कि $t \notin \Q(\sqrt2)$, कि $t + \sqrt2$
  अपरिमेय है, और यह कि कोई भी परिमेय $a, b$ $2^{1/3} = a + b\sqrt3$ संतुष्ट
  नहीं करते।
\end{enumerate}

\medskip \noindent\textbf{भाग VI --- सभी उपक्षेत्र, और संश्लेषण।}
\begin{enumerate}[resume]
  \item घातों की \omterm{def:b1:arith:divides}{विभाज्यता} सिद्ध कीजिए: यदि $\Q \subseteq A \subseteq B
  \subseteq \R$ ऐसे उपक्षेत्र हैं कि $\dim_\Q B$ परिमित है, तो $\dim_\Q A$
  $\dim_\Q B$ को \omterm{def:b1:arith:divides}{विभाजित करता है}। $M$ के उपक्षेत्रों की संभव विमाएँ क्या
  हैं?
  \item $\Q$ पर $2$ विमा वाले $M$ के \emph{सभी} उपक्षेत्र निर्धारित कीजिए।
  \emph{($w^2 \in \Q$ वाले $w = a + b\sqrt2 + c\sqrt3 + d\sqrt6 \in M
  \setminus \Q$ खोजने तक घटाइए: $w^2$ प्रसारित कीजिए और तीनों अपरिमेय
  \omterm{prop:b1:vspaces:coordinates}{निर्देशांक} शून्य कीजिए।)}
  \item $\dfrac1{1 + \sqrt2 + \sqrt3}$ को \omterm{def:b1:vspaces:free}{आधार} $(1, \sqrt2, \sqrt3, \sqrt6)$
  में व्यक्त कीजिए \emph{(सुचुने हुए संयुग्मियों से गुणा कीजिए)}।
  \item संश्लेषण, चार वाक्यों में: $\Q$-विमा वास्तविक संख्याओं के किसी
  \omterm{def:b1:logic:sets}{समुच्चय} का \emph{अंकगणितीय} अचर क्यों है; विमा की परिमितता हर अवयव के विषय
  में क्या बाध्य कर देती है (प्रश्न 13); तीन में से दो वाले नियम ने हवा में
  से प्रतिलोम कैसे पैदा किए (प्रश्न 12); और \omterm{pb:b1:findim:1}{मीनार नियम} किस बात को मना करता
  है (प्रश्न 20)। भाग III में सिद्ध प्रमेय का नाम बताइए।
\end{enumerate}
\end{problem}
